% ===================================================================
%  तक्ता क्रमांक ८ — कापणी.
%
%  Traduction, faite en 2026, en marathi d'aujourd'hui, sur l'ido de
%  texto/io. Regles de la colonne : voir 10-takta-01.tex. Deux
%  scenes, deux numerotations : les cles portent c1 et c2. Le
%  fac-simile compose « (13) » sans sa parenthese ouvrante ; la
%  traduction ecrit la reference bien formee, comme le fait deja la
%  colonne hindi.
%
%  LE FLEAU REVIENT, ET LE MARATHI FAIT EXACTEMENT CE QU'AVAIT FAIT
%  L'EGYPTIEN. Au tableau 9, la colonne egyptienne n'avait pas de mot
%  pour le fleau et avait ecrit عصي الدراس, « les batons de battage »,
%  parce que l'Egypte bat au نورج et non au fleau. Le marathi n'en a
%  pas davantage, pour la meme raison : le Maharashtra bat en faisant
%  pietiner les boeufs, ou a la machine. On ecrit donc मळणीच्या
%  काठ्या — mot pour mot la meme solution, dans une langue qui n'a
%  aucun rapport avec l'arabe. Deux colonnes, deux familles, un seul
%  trou, et il est dans la CHOSE, pas dans la langue.
%
%  TROIS EMPRUNTS ASSUMES, POUR LA MEME RAISON QUE LE MARRON CHAUD DU
%  TABLEAU 7 : le bleuet, la digitale, l'aubepine, le tilleul, le
%  chardonneret et la becasse ne sont pas des plantes ni des oiseaux
%  du Maharashtra. On les nomme par l'emprunt et l'on ne transpose
%  pas : ecrire un nom d'ici donnerait une phrase plus naturelle et
%  une traduction fausse. Le coquelicot, lui, tombe juste — खसखस
%  pousse en Inde depuis toujours.
%
%  ET LA MOISSON, ELLE, EST ENTIEREMENT MARATHE :
%
%    कुरण (9)     le paturage       सहाण (30)   la pierre a aiguiser
%    विळा (25)    la faux           कोयता (12)  la faucille
%    पेंढी (31)   la gerbe          मोळी (32)   le fagot
%    गंजी (44)    la meule          पेंढा       la paille
%    कणीस (29)    l'epi             झाडी (49)   le bois
%    झरा (20)     la source         ओढा (16)    le ruisseau
%    टेकडी (59)   la hauteur        सपाटी (39)  la plaine
%    चंडोल (41)   l'alouette        धोबी (19)   la bergeronnette
%    पाकोळी (56)  l'hirondelle      सरडा (2)    le lezard
%    घाणा (28)    le pressoir       घड          la grappe
%
%  DEUX NOMS POUR UN SEUL HOMME, LA OU L'IDO EN COMPOSE DEUX. Le
%  pastoro (10) est धनगर, le nom de la caste ; le muton-gardisto
%  (13) est मेंढपाळ, le nom de la fonction. L'ido devait composer ;
%  le marathi avait deja les deux mots, et ils ne veulent pas dire
%  la meme chose.
%
%  DEUX INVERSIONS PREVENUES au huitieme alinea — les culbutes sur
%  les tas de paille, et le champ fauche a la faux — MAIS SEPT
%  AUTRES ONT PASSE, et c'est renvoji.py qui les a vues : le berger
%  a la gourde, la bergeronnette a la source, le betail sous le
%  tilleul, l'alouette sur la plaine, le poisson qui happe
%  l'hamecon, les arbres sur la hauteur.
%
%  ET UN BLOC EN TROP, AU TROISIEME ALINEA DE LA SECONDE SCENE. Le
%  fac-simile coupe l'alinea au changement de feuillet et le reprend
%  sous « suite » ; j'en avais fait deux blocs, et renvoji.py a
%  repondu « cle inconnue ». Une traduction ne suit pas les
%  feuillets : les deux morceaux se rejoignent sous une seule cle.
% ===================================================================

%%K t08-noto-1 noto
\VUnotes{143.2mm}{%
\textit{(*) कारण हा शब्द नेमका नाही : ही जवसाची कापणी आहे, द्राक्षतोडणी आहे, की सफरचंदाची तोडणी?} \textit{Mondo XIV, पृष्ठ २४२ मध्ये सुचवलेला «} \VUgras{mesono} \textit{» हा शब्द वापरणे कदाचित बरे ठरेल. पण आतापर्यंत ती नवी धातू अकादमीने स्वीकारलेली नाही, आणि इडोच्या रूढ नियमांप्रमाणे ती चर्चेची वाट पाहत आहे.}%
}

%%K t08-apar-1 apar
\VUsaut{5.0mm}
\VUtitre{15.0pt}{60}{तक्ता क्रमांक ८}
\VUsaut{3.0mm}
\VUcentre{13.2pt}{90}{\textit{पहिला प्रवेश.}}
\VUsaut{3.0mm}
\VUpk{10.2pt}{40}{कापणी (*).}
\VUsaut{4.0mm}
\VUpk{10.2pt}{40}{शिवाराचे रूप.}
\VUsaut{3.0mm}

%%K t08-c1-01-1 p
१. --- \VUgras{उन्हाळ्याबरोबर} शिवाराचे रूप पुन्हा एकदा बदलले. तासात सोपवलेले बियाणे रुजले;
जमिनीतून हिरवे रोप बाहेर आले आणि त्याला कणसे धरली. दाणा भरला, मग तो पिकला, आणि आता तो कापून
घेण्याखेरीज काहीच उरले नाही.

%%K t08-c1-02-1 p
२. --- \VUgras{शिवारातली फुले} उमलली आहेत. \VUgras{कॉर्नफ्लॉवर} \textsuperscript{(1)},
\VUgras{खसखशीची फुले} \textsuperscript{(2)} आणि \VUgras{फॉक्सग्लोव्ह} \textsuperscript{(3)}
गव्हाच्या शेतांच्या एकसुरी रंगात आपल्या ताज्या \VUgras{पाकळ्यांचा} आनंदी विरोध मिसळतात.

%%K t08-c1-03-1 p
३. --- फळझाडांना पाने आली आहेत; फळ पिकले आहे. लहानशा \VUgras{चेरीच्या झाडाला}
\textsuperscript{(4)} सुंदर \VUgras{चेरी} \textsuperscript{(5)} लगडल्या आहेत, ज्या मुले मोठ्या
आनंदाने तोडतात आणि खातात.

%%K t08-c1-04-1 p
४. --- आतापासून जनावरे सगळा दिवस मोकळ्या हवेत घालवू शकतात.

%%K t08-c1-04-2 p
\VUgras{कोकरे} \textsuperscript{(6)} मोठी झाली आणि दाट लोकरीची सुंदर \VUgras{मेंढी} \textsuperscript{(7)} आणि सुंदर \VUgras{मेंढे} \textsuperscript{(8)} झाली. डेझीच्या फुलांनी नक्षीदार झालेल्या \VUgras{कुरणातले} \textsuperscript{(9)} \VUgras{गवत} ती शांतपणे चरतात.

%%K t08-c1-04-3 p
लवकरच \VUgras{ती} जनावरे कातरतील, म्हणजे त्यांची \VUgras{लोकर} मिळेल, जिच्यापासून आपल्या
कपड्यांचे कापड तयार होते.

%%K t08-tit-2 sub
\VUsaut{5.0mm}
\VUpk{10.2pt}{40}{धनगर.}
\VUsaut{3.4mm}

%%K t08-c1-05-1 p
५. --- \VUgras{धनगराचे} \textsuperscript{(10)} त्यांच्याकडे फारसे लक्ष दिसत नाही. क्षितिजावरून
जाणाऱ्या वादळाचीच त्याला काळजी आहे, आणि \VUgras{मेंढपाळ} \textsuperscript{(13)} तर आणखीच
बेफिकीर दिसतो — आपला \VUgras{तुंबा} \textsuperscript{(14)} ओठांजवळ नेणारा.

%%K t08-c1-06-1 p
६. --- ऊन असह्य आहे; कारण \VUgras{कुत्राही} \textsuperscript{(15)} तहान भागवायला येतो; तो
\VUgras{ओढ्याचे} \textsuperscript{(16)} थंड पाणी चाटतो आणि काठच्या गवतात दडलेल्या बिचाऱ्या
\VUgras{बेडकाला} \textsuperscript{(17)} घाबरवतो. तो स्वच्छ पाण्यात उडी मारतो, जिथे
\VUgras{कमळे} \textsuperscript{(18)} फुलली आहेत.

%%K t08-c1-07-1 p
७. --- एक \VUgras{धोबी} \textsuperscript{(19)} \VUgras{झऱ्याजवळ} \textsuperscript{(20)} अंघोळ
करतो — जो दोन खडकांमधून उसळतो — आणि \VUgras{काटेरी झुडपांखाली} \textsuperscript{(21)} व
\VUgras{हॉथॉर्नच्या झुडपांखाली} \textsuperscript{(22)} दडतो.

%%K t08-tit-3 sub
\VUsaut{5.0mm}
\VUpk{10.2pt}{40}{कापणी.}
\VUsaut{3.4mm}

%%K t08-c1-08-1 p
८. --- गावच्या मुलांसाठी गव्हाची कापणी हा फार मजेचा काळ असतो. ते आनंदाने \VUgras{कोलांट्या}
\textsuperscript{(24)} मारतात, \VUgras{पेंढ्याच्या ढिगांवर} \textsuperscript{(23)}. ते
\VUgras{कापणी करणाऱ्यांना} \textsuperscript{(26)} मदत करतात, आणि ते \VUgras{गव्हाचे शेत}
\textsuperscript{(27)} कापत असताना — आपल्या \VUgras{विळ्यांनी} \textsuperscript{(25)},
\VUgras{सहाणेवर} \textsuperscript{(30)} चांगली धार लावलेल्या — मुले गहू गोळा करतात आणि
त्याच्या \VUgras{पेंढ्या} \textsuperscript{(31)} किंवा \VUgras{मोळ्या} \textsuperscript{(32)}
बांधतात.

%%K t08-c1-09-1 p
९. --- कधी कधी त्यांच्यातल्या सर्वात मोठ्याला \VUgras{कोयत्याने} \textsuperscript{(12)} ते
\VUgras{सोनेरी दांडे} \textsuperscript{(28)} कापू देतात, जे दाण्यांनी भरलेली जड \VUgras{कणसे}
\textsuperscript{(29)} पेलतात; आणि लहानगे \VUgras{जनावरांवर} \textsuperscript{(33)} नजर ठेवतात
— जी \VUgras{कुरणात} \textsuperscript{(34)} चरतात, मोठ्या \VUgras{लिंडेनच्या झाडाच्या}
\textsuperscript{(35)} सावलीत — आणि आपल्या \VUgras{जाळ्याने} \textsuperscript{(36)} सुंदर
\VUgras{फुलपाखरे} पकडतात.

%%K t08-c1-09-2 p
काही जण तर \VUgras{गाडीवर} \textsuperscript{(37)} चढून, \VUgras{काट्याने}
\textsuperscript{(38)} त्यांच्याकडे दिल्या जाणाऱ्या पेंढ्या रचतात.

%%K t08-c1-10-1 p
१०. --- सगळ्या \VUgras{सपाटीवर} \textsuperscript{(39)} चैतन्य पसरले आहे, जिथे \VUgras{चंडोल}
\textsuperscript{(41)} उडू लागतात. सगळे कापणीत गुंतले आहेत. पण गरीब लोक जुनीच अवजारे वापरत
राहतात — \VUgras{विळे, कोयते} आणि \VUgras{मळणीच्या काठ्या} \textsuperscript{(40)} —, तर
श्रीमंत जमीनदार आपला गहू किंवा आपले \VUgras{राय} \textsuperscript{(43)}
\VUgras{कापणीच्या यंत्राने} \textsuperscript{(42)} कापून घेतात. मग पेंढ्या खळ्यात न नेता,
तिथेच मोठ्या \VUgras{गंजी} \textsuperscript{(44)} रचतात.

%%K t08-c1-11-1 p
११. --- मग \VUgras{मळणी यंत्र} \textsuperscript{(47)} आणतात, जे \VUgras{वाफेचे इंजिन}
\textsuperscript{(45)} \VUgras{पट्ट्याने} \textsuperscript{(46)} चालवते, आणि अशा रीतीने दाणा
\VUgras{पेंढ्यापासून} वेगळा होतो — ज्या शेतात तो पेरला होता ते न सोडता.

%%K t08-tit-4 sub
\VUsaut{5.0mm}
\VUpk{10.2pt}{40}{वादळ.}
\VUsaut{3.4mm}

%%K t08-c1-12-1 p
१२. --- कापणीच्या काळात अनेकदा \VUgras{मोठी वादळे} \textsuperscript{(53)} होतात. आधी दुरून
\VUgras{ढगांचा गडगडाट} ऐकू येतो; \VUgras{पश्चिमेचे वादळी वारे} \textsuperscript{(54)} सुटतात
आणि धापा टाकत \VUgras{झाडीतली} \textsuperscript{(49)} सर्वात जाड झाडे वाकवतात. काळे
\VUgras{ढग} \textsuperscript{(56)} आकाशावर चाल करतात; \VUgras{विजेच्या} \textsuperscript{(55)}
डोळे दिपवणाऱ्या नागमोडी रेषा त्यांना चिरत जातात. \VUgras{पाऊस} आणि कधी कधी \VUgras{गारा} पडू
लागतात, कापणीला तयार असलेले पीक आडवे करत.

%%K t08-c1-13-1 p
१३. --- वीज अनेकदा एखाद्या इमारतीला आग लावते. अशीच गेल्या वर्षी \VUgras{पवनचक्कीचे}
\textsuperscript{(50)} छप्पर जळून राख झाले आणि तिची दोन \VUgras{पाती} \textsuperscript{(51)}
मोडली.

%%K t08-c1-14-1 p
१४. --- काही तासांनी पुन्हा शांतता पसरते. क्षितिजावर तेजस्वी \VUgras{इंद्रधनुष्य}
\textsuperscript{(52)} उमटते आणि काही क्षण थांबलेले निसर्गाचे जीवन पुन्हा सुरू होते.
\VUgras{झोपड्यांमध्ये} \textsuperscript{(48)} आश्रय घेतलेले गावकरी पुन्हा कामाला लागतात आणि
नुकत्याच सोसलेल्या नुकसानाची भरपाई करायचा धीराने प्रयत्न करतात.

%%K t08-tit-5 sub
\VUsaut{6.0mm}
\VUcentre{13.2pt}{90}{\textit{दुसरा प्रवेश.}}
\VUsaut{3.6mm}

%%K t08-tit-6 sub
\VUpk{10.2pt}{40}{शिकार. --- द्राक्षतोडणी.}
\VUsaut{1.4mm}
\VUpk{10.2pt}{40}{मासेमारी.}
\VUsaut{4.0mm}

%%K t08-c2-01-1 p
१. --- प्रदेशाप्रमाणे \VUgras{ऑगस्टच्या} शेवटी किंवा \VUgras{सप्टेंबरच्या} मध्यावर शिकारीचा
हंगाम सुरू होतो. सूर्याच्या पहिल्या किरणांनी \VUgras{सरडे} \textsuperscript{(2)} आपल्या
बिळातून बाहेर काढले न काढले तोच \VUgras{शिकारी} \textsuperscript{(5)} आपल्या
\VUgras{कुत्र्यांसह} \textsuperscript{(4)} वाटेला लागतो.

%%K t08-c2-02-1 p
२. --- त्याने जाड कापडाचा आपला मजबूत \VUgras{शिकारीचा पोशाख} \textsuperscript{(9)} चढवला,
पायांत कातडी \VUgras{घोटेबंद} घातले — ज्यांच्या जोरावर तो त्या ओल्या गवतातून चालू शकेल, जिथे
\VUgras{मोठे बेडूक} \textsuperscript{(1)} दडतात; त्याने आपली \VUgras{बंदूक}
\textsuperscript{(6)} आणि आपली \VUgras{शिकारीची पिशवी} \textsuperscript{(10)} घेतली, आणि हे बघ
तो निघाला.

%%K t08-c2-03-1 p
३. --- पाठलागाच्या नादात तो द्राक्षमळ्याच्या मधोमध पोहोचतो, जिथे द्राक्षतोडणी जोरात चालली आहे.
द्राक्षमळेवाल्याची मुले, जी एरवी वाटेवर शेळ्या राखतात, आज त्यांना वाटेल तिथे चरू देतात, आणि
आपले स्वातंत्र्य वापरून पळून जायला कमी करणार नाही अशा म्हाताऱ्या \VUgras{बोकडाला}
\textsuperscript{(3)} खुंटाला बांधून, तीही आपल्या परीने या मौजेत वाटा उचलतात. एक जण
\VUgras{तोडणीच्या टोपलीतला} \textsuperscript{(12)} द्राक्षांचा घड उचलतो आणि \VUgras{रस}
\textsuperscript{(14)} एका \VUgras{पेल्यात} \textsuperscript{(13)} गाळत नवा रस (न आंबवलेली
दारू) करण्याची मौज करतो. दुसरा नुकत्याच विणलेल्या \VUgras{पानांच्या मुकुटाने}
\textsuperscript{(15)} आपले केस सजवतो. त्यांच्याजवळ निर्लज्ज \VUgras{चिमण्या}
\textsuperscript{(16)} द्राक्षे चोरायला घाण्यासमोरपर्यंत येतात.


%%K t08-c2-04-1 p
४. --- मुले मौज करत असताना पुरुष काम करतात. \VUgras{द्राक्षतोडणारे} \textsuperscript{(18)}
\VUgras{पाठीवरच्या टोपल्यांनी} \textsuperscript{(17)} द्राक्षे तळघरात नेतात;
\VUgras{पिंपे बनवणारा} \textsuperscript{(19)} \VUgras{पिंपे} \textsuperscript{(20)} दुरुस्त
करतो, \VUgras{फळ्या} \textsuperscript{(22)} आणि \VUgras{तळ} \textsuperscript{(21)} चांगल्या
स्थितीत आहेत का ते तपासतो, आणि तुटलेल्या \VUgras{लोखंडी कड्या} \textsuperscript{(23)} बदलतो.
त्यानेच \VUgras{घंगाळीही} \textsuperscript{(25)} तयार केली, ज्यांत \VUgras{द्राक्षे}
\textsuperscript{(26)} आंबवायला ओततात.

%%K t08-c2-05-1 p
५. --- आंबणे पुरे झाले की दारू काढून घेतात आणि उरलेला चोथा \VUgras{घाण्यावर}
\textsuperscript{(28)} ठेवतात, आणि \VUgras{मोठ्या स्क्रूने} \textsuperscript{(29)} हळूहळू दाबत
रस पिळून काढतात, जो \VUgras{कुंडात} \textsuperscript{(32)} वाहतो, आणि त्यातूनही \VUgras{वाइन}
\textsuperscript{(31)} मिळते.

%%K t08-tit-8 sub
\VUsaut{6.0mm}
\VUpk{10.2pt}{40}{बिअर आणि सायडर.}
\VUsaut{3.4mm}

%%K t08-c2-06-1 p
६. --- \VUgras{तळघराच्या} \textsuperscript{(27)} भिंतीवर \VUgras{हॉप्सची}
\textsuperscript{(30)} एक फांदी चढते. तिथे ती केवळ शोभेची वेल आहे, पण काही देशांत, जिथे
द्राक्षवेल फार दुर्मीळ आहे, तिथे \VUgras{बिअर} करण्यासाठी हॉप्स लावतात.

%%K t08-c2-07-1 p
७. --- लहानशा \VUgras{सफरचंदाच्या झाडाला} \textsuperscript{(33)} सफरचंदे लगडली आहेत, जी
पिकल्यावर उत्तम \VUgras{सायडर} देतील. आपल्या \VUgras{पिंजऱ्यात} \textsuperscript{(34)}
\VUgras{कॅनरी} \textsuperscript{(35)} आणि \VUgras{गोल्डफिंच} \textsuperscript{(36)} मोकळ्या
बागडणाऱ्या चिमण्यांकडे मत्सराच्या नजरेने पाहतात.

%%K t08-c2-08-1 p
८. --- नजर पोहोचेल तिथवर, टेकड्यांच्या उतारावर \VUgras{वेलींचे खांब} \textsuperscript{(40)}
रांगेत उभे आहेत, जे जड \VUgras{घडांनी} लगडलेल्या देखण्या \VUgras{वेली} \textsuperscript{(39)}
पेलतात.

%%K t08-c2-08-2 p
वर्षभर \VUgras{द्राक्षमळेवाल्याने} \textsuperscript{(42)} आपल्या \VUgras{द्राक्षमळ्याची}
\textsuperscript{(38)} काळजी घ्यायला कष्ट केले, आणि आता सुंदर पिकाने त्याच्या श्रमाचे फळ
त्याला मिळाले आहे. \VUgras{द्राक्षे तोडणाऱ्या बायका} \textsuperscript{(41)} \VUgras{खेचरांनी}
\textsuperscript{(43)} ओढलेल्या गाडीवर ठेवलेली घंगाळी भरतात, आणि सगळे आनंदात आहेत.

%%K t08-tit-9 sub
\VUsaut{5.0mm}
\VUpk{10.2pt}{40}{मासेमारी.}
\VUsaut{3.4mm}

%%K t08-c2-09-1 p
९. --- तसेच \VUgras{गळ टाकणाऱ्यांचेही} \textsuperscript{(44)}, जे सकाळपासून आपल्या
\VUgras{गळाच्या काठ्या} \textsuperscript{(45)} घेऊन पाण्याच्या काठावर येऊन बसले आहेत.

%%K t08-c2-09-2 p
\VUgras{मासा} \textsuperscript{(47)} \VUgras{गळ} तोंडात घेतो, ते आपला \VUgras{दोरा}
\textsuperscript{(46)} पाण्यात सोडताच; आणि \VUgras{आमिष} पुन्हा लावायला त्यांना उसंत मिळते न
मिळते तोच नवा बळी दोऱ्याच्या टोकाला तडफडू लागतो.

%%K t08-c2-09-3 p
तरीही \VUgras{जाळे टाकणारा} \textsuperscript{(48)}, जो आपल्या \VUgras{होडीतून}
\textsuperscript{(49)} \VUgras{नदीच्या} \textsuperscript{(50)} मधोमध \VUgras{जाळे} फेकतो, तो
जास्त मासे पकडतो.

%%K t08-c2-10-1 p
१०. --- शरद ऋतू हा चांगल्या फेरफटक्यांचा हंगाम : पायी, \VUgras{घोड्यावरून}
\textsuperscript{(52)}, \VUgras{सायकलवरून} \textsuperscript{(53)}, \VUgras{मोटारीने}
\textsuperscript{(54)}. \VUgras{रस्ते} \textsuperscript{(51)} स्वच्छ आहेत आणि शेवटच्या
चांगल्या दिवसांचा लोक फायदा घेतात, कारण हे बघ \VUgras{पाकोळ्या} \textsuperscript{(56)} आधीच
छपरांवर जमू लागल्या आहेत, अधिक उबदार देशांकडे पंख पसरून उडून जाण्यासाठी. म्हाताऱ्या
\VUgras{अक्रोडाच्या झाडाची} \textsuperscript{(57)} पाने पिवळी पडत आहेत, \VUgras{अक्रोड} गळत
आहेत, आणि \VUgras{पुष्पगुच्छासारखी} \textsuperscript{(58)} गटागटाने उभी असलेली उंच
\VUgras{झाडे} — \VUgras{टेकडीवरची} \textsuperscript{(59)} — लवकरच निष्पर्ण होतील.

%%K t08-c2-11-1 p
११. --- \VUgras{सूर्य} \textsuperscript{(60)} उशिरा उगवतो, दिवस लहान होत जातात, \VUgras{सकाळी}
बऱ्यापैकी बोचरी थंडी जाणवू लागते, आणि हे बघ \VUgras{वुडकॉकचे} \textsuperscript{(61)} थवे आधीच
आपले पाहुणचार न करणारे हवामान सोडून चालले आहेत.

\VUsaut{6.0mm}
\VUfilet{22mm}
